\section{Materiales}

\subsection{Lista de Materiales de la Práctica}
En la Tabla~\ref{tab:materiales} se desglosan los materiales solicitados para la práctica y los componentes equivalentes o adicionales implementados en el proyecto.

\begin{table}[H]
\centering
\small
\begin{tabularx}{\textwidth}{c p{3.2cm} p{4.0cm} X c}
\toprule
\textbf{\#} & \textbf{Requerido en Guía} & \textbf{Componente Utilizado} & \textbf{Función y Justificación} & \textbf{Cant.} \\
\midrule
1 & 1 ESP32 & ESP32 DevKit V1 (30 pines) & Microcontrolador con Wi-Fi, cliente MQTT y GPIO. & 1 \\
2 & 1 Sensor DHT11 & Sensor digital GY-SHT41 \newline \textit{(Sensirion SHT41, I²C)} & Mide temperatura ($\pm 0.2\,^\circ\text{C}$) y humedad ($\pm 1.8\%$). Supera al DHT11 y usa bus estándar I²C. & 1 \\
3 & 1 LED & LED difuso de 5 mm rojo & Actuador conmutable conectado a GPIO 2. & 1 \\
4 & Resistencias & Resistor $220\,\Omega$, $1/4\,\text{W}$ & Limitador de corriente de polarización del LED. & 1 \\
5 & Protoboard & Protoboard de 830 puntos & Placa de pruebas para el montaje sin soldadura. & 1 \\
6 & Cables & Jumpers macho-macho & Conexiones de bus I²C, GPIO, 3.3V y GND. & $\sim 12$ \\
7 & Computadora e Internet & Laptop con Linux (hotspot) & Alojamiento de Mosquitto broker y Node-RED. & 1 \\
\midrule
\multicolumn{5}{l}{\proyectoBadge \textbf{Componentes adicionales incorporados en el proyecto:}} \\
8 & Pantalla local & OLED 128$\times$32 SSD1306 (I²C) & Muestra temp/humedad y estado del LED in situ. & 1 \\
\bottomrule
\end{tabularx}
\caption{Lista de materiales requeridos y componentes utilizados.}
\label{tab:materiales}
\end{table}

\subsection{Distinción Técnica: DHT11 frente a SHT41 y Diagnóstico BME680}
\begin{itemize}
    \item \practicaBadge \textbf{Sensor de Temperatura y Humedad:} La guía solicita un DHT11 (protocolo propietario 1-wire, resolución de 8 bits, precisión $\pm 2^\circ\text{C}$ / $\pm 5\%$). En este trabajo se optó por el sensor digital \textbf{SHT41} de Sensirion por su protocolo estándar I²C (dirección \texttt{0x44}), resolución de 16 bits ($\pm 0.2^\circ\text{C}$ / $\pm 1.8\%$) y operación no bloqueante en el microcontrolador.
    \item \proyectoBadge \textbf{Antecedente del BME680:} Inicialmente se probó un sensor BME680 (temperatura, humedad, presión y gas). Durante las pruebas diagnósticas se constató que su registro interno de temperatura arrojaba invariablemente \texttt{adc\_T = 0} por falla de fabricación en el silicio, por lo que fue sustituido por el SHT41, focalizando el proyecto en confort térmico.
\end{itemize}
